
\usepackage[T2A]{fontenc}
\usepackage[utf8]{inputenc}
\usepackage[russian]{babel}
\usepackage{amsthm}
\usepackage{hyperref}
\usepackage{amsmath,amssymb}
\usepackage{booktabs}
\usepackage{longtable}
\usepackage{geometry}
\geometry{margin=1in}

\section*{Связанная термомеханическая модель с учётом продольных и поперечных волн}

В данной работе используется система уравнений термоупругости, в которой 
коэффициенты материала зависят от давления. 
Экспериментальные данные по плотности и скоростям распространения упругих волн 
берутся из таблицы свойств пород при высоких давлениях. 
Эти данные позволяют выразить модули упругости и динамические характеристики среды 
как функции давления глубины.

\subsection*{1. Уравнение теплопроводности}

\begin{equation}
\rho(P)\, c_p(T)\, \frac{\partial T}{\partial t}
=
\nabla \cdot \left( k(T)\, \nabla T \right)
+ Q,
\label{eq:heat}
\end{equation}

где $T$ — температура,
$\rho(P)$ — плотность породы из экспериментальной таблицы,
$c_p$ — теплоёмкость,
$k$ — теплопроводность,
$Q$ — объёмные тепловые источники.

\subsection*{2. Уравнения движения}

Уравнение импульса для упругой среды с учётом зависимости плотности от давления имеет вид:
\begin{equation}
\rho(P)\, \frac{\partial^2 u_i}{\partial t^2}
=
\nabla_j \sigma_{ij}
+ f_i,
\label{eq:motion}
\end{equation}

где $u_i$ — компоненты вектора смещений,
$f_i$ — массовые силы,
$\sigma_{ij}$ — тензор напряжений.

\subsection*{3. Линейная связь между напряжениями и деформациями}

Для изотропной упругой среды:

\begin{equation}
\sigma_{ij}
=
\lambda(P)\, \delta_{ij}\, \varepsilon_{kk}
+
2 G(P)\, \varepsilon_{ij},
\label{eq:constitutive}
\end{equation}

где $\lambda(P)$ и $G(P)$ — ляме-параметры,
зависящие от давления,
\[
\varepsilon_{ij}
=
\frac{1}{2}
\left(
\frac{\partial u_i}{\partial x_j}
+
\frac{\partial u_j}{\partial x_i}
\right)
\]
— тензор деформаций.

\subsection*{4. Выражение модулей упругости через скорости волн}

Экспериментальная таблица содержит скорости продольных (компрессионных) и поперечных волн.
Эти величины используются для определения модулей среды.

\paragraph{Поперечные волны:}
\begin{equation}
V_s(P)
=
\sqrt{\frac{G(P)}{\rho(P)}}.
\label{eq:Vs}
\end{equation}

\paragraph{Продольные волны:}
\begin{equation}
V_p(P)
=
\sqrt{
\frac{
K(P) + \frac{4}{3} G(P)
}{
\rho(P)
}
}.
\label{eq:Vp}
\end{equation}

\subsection*{5. Определение модулей $K(P)$ и $G(P)$}

Из уравнений \eqref{eq:Vs}–\eqref{eq:Vp} получаем:

\begin{equation}
G(P) = \rho(P)\, V_s^2(P),
\label{eq:G}
\end{equation}

\begin{equation}
K(P) = \rho(P)\, V_p^2(P) - \frac{4}{3} G(P).
\label{eq:K}
\end{equation}

\subsection*{6. Система, используемая в численном моделировании}

Полная связанная система:

\[
\begin{cases}
\rho(P)\, c_p(T)\, T_t - \nabla \cdot (k \nabla T) = Q,\\[6pt]
\rho(P)\, u_{tt} - \nabla \cdot \sigma = f,\\[6pt]
\sigma = C(P,T) : \varepsilon,\\[6pt]
G(P) = \rho(P)\, V_s^2(P),\\[4pt]
K(P) = \rho(P)\, V_p^2(P) - \dfrac{4}{3} G(P).\\
\end{cases}
\]

Здесь функции $\rho(P)$, $V_p(P)$ и $V_s(P)$ задаются 
экспериментальными таблицами для пород различного состава.
Материальные параметры обновляются на каждом временном шаге,
что обеспечивает учёт реального поведения среды при увеличении давления.
